\chapter{Sequences and Their Limits}

For most of this book a limit has been a process \emph{along a continuum}. You
let $x$ slide smoothly toward some point $a$ and you watched what $f(x)$ did. The
input could come from the left, from the right, through rationals, through
irrationals -- any way at all, as long as it crept up on $a$. That freedom was
the whole content of the epsilon-delta definition.

Now we change the shape of the approach. Instead of a continuous variable
gliding toward a point, picture a discrete \emph{list} of numbers, one after
another, marching out forever:
\[
  a_1,\; a_2,\; a_3,\; a_4,\; \dots,\; a_n,\; \dots
\]
There is a first term, a second term, a thousandth term, but never a last term.
The question is no longer ``what happens as $x \to a$?'' but ``where is this list
\emph{headed} as we walk down it?'' Does it settle on a single number? Does it
wander forever? Does it blow up?

This is the idea of a \emph{sequence} and its \emph{limit}, and it is the gateway
to the second half of calculus. Every infinite series -- every Taylor expansion,
every $\pi = 4(1 - \tfrac13 + \tfrac15 - \cdots)$, every decimal that never
terminates -- is secretly a sequence in disguise: the sequence of its partial
sums. Newton's method produces a sequence of guesses. The trapezoid rule
produces a sequence of approximations. Even a real number like $\sqrt{2}$ is, at
bottom, the destination of a sequence of rationals. Understand how a list
of numbers can \emph{converge} and all of these come within reach at once.

The chapter proceeds in three stages. We begin by pinning down what a sequence is
and what it means for one to have a limit -- the same challenge-and-response game
you already know, re-tuned for the integers. From there we assemble the working
tools: the squeeze theorem for sequences, a short table of limits you will use
constantly, and the deep fact that a sequence which climbs but stays bounded
\emph{must} converge. Last comes recursion -- sequences in which each term is
built from the one before -- together with the idea of a \emph{fixed
point}, which is where such sequences come to rest.

\chapterline

\section{What a sequence really is}
\label{ch28:what-is-a-sequence}

Strip away the dots and a sequence is something you already understand: a
function. The only twist is its domain.

\begin{definition}[Sequence]
\label{ch28:def-sequence}
A \emph{sequence} of real numbers is a function $a \colon \N \to \R$ whose
domain is the natural numbers $\N = \{1, 2, 3, \dots\}$. Instead of writing
$a(n)$ we write $a_n$ and call it the \emph{$n$th term}. We denote the whole
sequence by $(a_n)_{n=1}^{\infty}$, or just $(a_n)$.
\end{definition}

That is the entire definition. A sequence is a function whose inputs are
$1, 2, 3, \dots$ rather than a continuous interval. Where an ordinary function
$f$ has a smooth graph you can trace with a pencil, a sequence has a graph made
of \emph{isolated dots}, one above each integer on the horizontal axis.

\begin{intuitionbox}
Three ways to read the same object, all useful:
\begin{itemize}
  \item \textbf{As a list.} $\tfrac11, \tfrac12, \tfrac13, \tfrac14, \dots$ --
        the thing you write down term by term.
  \item \textbf{As a function.} $a_n = \tfrac1n$, the rule that turns an input
        $n$ into an output. Plug in $n=7$, read off $\tfrac17$.
  \item \textbf{As a picture.} A scatter of dots at heights
        $1, \tfrac12, \tfrac13, \dots$ above the integers $1, 2, 3, \dots$,
        sinking toward the horizontal axis but never touching it.
\end{itemize}
Keep all three in mind. The list view is how you compute, the function view is
how you reason, and the picture view is how you \emph{see} where the sequence is
going.
\end{intuitionbox}

There are three common ways to hand someone a sequence.

\textbf{By a closed formula.} You give $a_n$ as an explicit expression in $n$.
For example $a_n = \dfrac{n}{n+1}$ produces $\tfrac12, \tfrac23, \tfrac34,
\tfrac45, \dots$. To find the thousandth term you just substitute: $a_{1000} =
\tfrac{1000}{1001}$. No bookkeeping required.

\textbf{By recursion.} You give the first term (or first few) and a rule that
builds each new term from earlier ones. The classic is the Fibonacci sequence,
$F_1 = F_2 = 1$ and $F_{n} = F_{n-1} + F_{n-2}$, which marches out as $1, 1, 2,
3, 5, 8, 13, \dots$. Recursion is powerful but slippery: to know $a_{1000}$ you
must, in principle, walk through all $999$ terms before it. Section
\ref{ch28:recursion} is devoted to taming such sequences.

\textbf{By description.} Sometimes the rule is verbal: ``let $p_n$ be the $n$th
prime,'' giving $2, 3, 5, 7, 11, \dots$. There is no tidy formula, but it is a
perfectly good function $\N \to \R$.

\begin{walkbox}
Write out the first five terms of $a_n = \dfrac{(-1)^n}{n}$. What is the sign
doing, and what are the magnitudes doing? Where do you think this list is
headed?

\smallskip
The terms are $-1,\ \tfrac12,\ -\tfrac13,\ \tfrac14,\ -\tfrac15$. The sign
alternates; the magnitudes $1, \tfrac12, \tfrac13, \tfrac14, \tfrac15$ shrink
toward zero. So the dots bounce above and below the axis but the bounces get
ever smaller. The list is being squeezed toward $0$ from both sides. Hold that
image -- it is exactly the case the squeeze theorem (Section \ref{ch28:squeeze})
will handle.
\end{walkbox}

\begin{marginnotebox}
A word on indexing. Some sequences naturally start at $n=0$ (powers, $2^0,
2^1, 2^2, \dots$) and some at $n=1$. Nothing about a limit depends on where the
list starts -- changing or deleting the first million terms cannot change where
the \emph{tail} is headed. We will start at $n=1$ unless a formula begs to start
at $0$, and we will be relaxed about it, because the mathematics is relaxed about
it.
\end{marginnotebox}

\section{Convergence: the destination of a list}
\label{ch28:convergence}

Look again at $a_n = \tfrac1n$. The terms are $1, \tfrac12, \tfrac13, \tfrac14,
\dots$. Nobody hesitates to say this sequence ``goes to $0$.'' But what,
\emph{precisely}, does that claim assert? Not that any term equals zero -- none
does. Not merely that the terms decrease -- the sequence $1, \tfrac12, \tfrac12,
\tfrac12, \dots$ also decreases (weakly) but goes to $\tfrac12$, not $0$. The
real content is this:

\begin{quote}
\emph{No matter how small a margin of error you name, eventually -- from some
term onward -- every term of the sequence lands inside that margin of $0$.}
\end{quote}

This is the challenge-and-response game again, the same one that defined
$\lim_{x\to a} f(x)$, but cleaner, because there is no $\delta$ to find. The
input does not slide toward a point; it just \emph{increases}. So the response
to a challenge is not ``get within $\delta$ of $a$'' but simply ``go far enough
out in the list.''

\begin{intuitionbox}
\textbf{The game.} A skeptic challenges you with a tiny positive number
$\eps$ -- a tolerance band of half-width $\eps$ drawn around the candidate
limit $L$. Your job is to answer with a cutoff $N$: a place in the list beyond
which \emph{all} terms lie inside the band. If you can meet \emph{every}
challenge, no matter how small the $\eps$, then $L$ is the limit. If there is
even one $\eps$ you cannot answer, $L$ is not the limit.

The skeptic moves first and can make $\eps$ as cruelly small as desired. You
move second, so you may choose $N$ \emph{after} seeing $\eps$ -- and $N$ is
allowed to be enormous. Small $\eps$ usually forces large $N$; that is fine. The
only thing that matters is that for each $\eps$ \emph{some} workable $N$ exists.
\end{intuitionbox}

\begin{definition}[Limit of a sequence]
\label{ch28:def-limit}
A sequence $(a_n)$ \emph{converges} to the number $L$, written
\[
  \lim_{n\to\infty} a_n = L
  \qquad\text{or}\qquad
  a_n \To L,
\]
if for every $\eps > 0$ there exists a natural number $N$ such that
\[
  n \ge N \quad\Longrightarrow\quad \abs{a_n - L} < \eps .
\]
If no such $L$ exists, the sequence \emph{diverges}.
\end{definition}

Read that implication out loud as a sentence in the game: ``once you go past the
cutoff $N$, every term is within $\eps$ of $L$.'' The quantity $\abs{a_n - L}$ is
the distance from the $n$th term to the candidate limit; demanding it be less
than $\eps$ is demanding the term sit inside the tolerance band.

\begin{figure}[H]
\centering
\begin{tikzpicture}
\begin{axis}[
  width=12cm, height=6.4cm,
  xlabel={$n$}, ylabel={$a_n$},
  xmin=0, xmax=21, ymin=-0.15, ymax=1.15,
  axis lines=left,
  xtick={1,5,10,15,20},
  ytick={0,0.5,1},
  clip=false,
]
% tolerance band around L=0
\addplot[draw=none, name path=top] coordinates {(0,0.12)(21,0.12)};
\addplot[draw=none, name path=bot] coordinates {(0,-0.12)(21,-0.12)};
\addplot[blue!12] fill between[of=top and bot];
\addplot[densely dashed, blue!55!black] coordinates {(0,0)(21,0)};
\node[blue!55!black, anchor=west] at (axis cs:14.5,0.12) {\small $L+\eps$};
\node[blue!55!black, anchor=west] at (axis cs:14.5,-0.12) {\small $L-\eps$};
% the sequence 1/n
\addplot[only marks, mark=*, mark size=1.6pt, curveblue]
  coordinates {(1,1)(2,0.5)(3,0.3333)(4,0.25)(5,0.2)(6,0.1667)(7,0.1429)
  (8,0.125)(9,0.1111)(10,0.1)(11,0.0909)(12,0.0833)(13,0.0769)(14,0.0714)
  (15,0.0667)(16,0.0625)(17,0.0588)(18,0.0556)(19,0.0526)(20,0.05)};
% mark the cutoff N
\draw[trapred, thick] (axis cs:9,-0.15) -- (axis cs:9,1.15);
\node[trapred, anchor=south west] at (axis cs:9,0.78) {\small $N$};
\end{axis}
\end{tikzpicture}
\caption{The sequence $a_n = 1/n$ converging to $L=0$. The shaded strip is the
tolerance band $\abs{a_n - 0} < \eps$ for $\eps = 0.12$. To the right of the red
cutoff $N$ (here $N=9$) every dot lies inside the band. Shrink $\eps$ and the
band narrows, pushing the cutoff $N$ rightward -- but a workable $N$ always
exists.}
\label{ch28:fig-oneovern}
\end{figure}

\begin{example}[Proving $1/n \to 0$ from the definition]
\label{ch28:ex-oneovern}
We claim $\lim_{n\to\infty}\tfrac1n = 0$. Let $\eps > 0$ be given by the
skeptic. We must produce $N$ so that $n \ge N$ forces $\abs{\tfrac1n - 0} =
\tfrac1n < \eps$.

Solve the target inequality for $n$: $\tfrac1n < \eps$ exactly when $n >
\tfrac1\eps$. So \emph{any} integer $N$ larger than $\tfrac1\eps$ will do; the
existence of such an integer is guaranteed by the Archimedean property of the
real numbers (no real number is larger than every integer). Concretely, take
$N = \floor{1/\eps} + 1$.

Now verify. If $n \ge N > \tfrac1\eps$, then $n > \tfrac1\eps$, hence $\tfrac1n
< \eps$, hence $\abs{\tfrac1n - 0} < \eps$. The challenge is met, and since
$\eps$ was arbitrary, $\tfrac1n \to 0$. $\qed$
\end{example}

\begin{whybox}
Notice the \emph{shape} of that proof, because every epsilon-$N$ proof has it.
\begin{enumerate}
  \item \textbf{Let $\eps > 0$ be given.} You do not get to choose it; you must
        handle all of them at once by handling a typical one.
  \item \textbf{Work backwards.} Start from the goal $\abs{a_n - L} < \eps$ and
        solve for what $n$ must satisfy. This is scratch work -- it tells you how
        big $N$ has to be.
  \item \textbf{Declare $N$.} Announce the cutoff your scratch work revealed.
  \item \textbf{Verify forwards.} Assume $n \ge N$ and chain the inequalities to
        reach $\abs{a_n - L} < \eps$. This is the actual proof; steps 2--3 were
        the search.
\end{enumerate}
The search and the proof run in opposite directions. Beginners blur them
together and get confused about which inequality they are allowed to assume.
Keep them separate: solve backward in the margin, then write forward on the page.
\end{whybox}

\subsection*{Divergence, and a sequence with no limit}

A sequence diverges if no $L$ works. There are two flavors worth naming.

\textbf{Divergence to infinity.} The terms $a_n = n$ grow without bound. They
cannot converge to any finite $L$, because once the terms exceed $L + 1$ they can
never return to within $\eps = 1$ of $L$. We write $a_n \to \infty$ as shorthand,
but be clear: $\infty$ is not a number the sequence reaches; it is a description
of unbounded growth.

\textbf{Oscillation.} The terms $a_n = (-1)^n$ cycle $-1, 1, -1, 1, \dots$
forever. This sequence is bounded -- it never escapes $[-1, 1]$ -- yet it has no
limit. If we proposed $L = 1$, the term $-1$ keeps appearing arbitrarily far out,
so the choice $\eps = 1$ is never met from any cutoff on. The same defeat befalls
$L = -1$, or any $L$ in between. No single destination; no limit.

\begin{trapbox}
\textbf{Bounded does not mean convergent.} The sequence $(-1)^n$ stays trapped in
$[-1,1]$ but never settles. Convergence requires the terms to \emph{pile up} near
one number, not merely to stay in a finite range. We will see in Section
\ref{ch28:monotone} the extra ingredient -- \emph{monotonicity} -- that upgrades
``bounded'' into ``convergent.''
\end{trapbox}

\begin{historybox}
The notation $\lim$ comes from the Latin \emph{limes}, a boundary or frontier --
the Romans used it for the fortified edge of empire. The careful epsilon-$N$
formulation is the work of the nineteenth century, chiefly Bernard Bolzano,
Augustin-Louis Cauchy, and above all Karl Weierstrass, whose Berlin lectures in
the 1860s hammered the $\eps$ and the $N$ into the form you just read.

Before them, ``$a_n$ approaches $L$'' was a verbal picture, and pictures led good
mathematicians astray. Cauchy himself once ``proved'' that an infinite sum of
continuous functions is continuous -- a statement that is simply false, as
Fourier's heat-conduction series showed by summing smooth waves into a sawtooth
with a jump. The fix was not better intuition but a definition sharp enough to
catch the error. That is the whole moral of the epsilon-$N$ machinery: it lets
you be \emph{right} about infinity, not just confident.
\end{historybox}

\section{The limit laws, and why they transfer}
\label{ch28:limit-laws}

You spent a chapter proving that limits of functions respect addition,
multiplication, and division. The good news: sequence limits inherit every one
of those laws, with proofs that are if anything simpler, because there is no
$\delta$ to juggle.

\begin{theorem}[Algebra of sequence limits]
\label{ch28:thm-laws}
Suppose $a_n \to A$ and $b_n \to B$. Then
\[
  a_n + b_n \To A + B, \qquad
  a_n\,b_n \To A\,B, \qquad
  c\,a_n \To c\,A \ \ (\text{constant } c),
\]
and, provided $B \ne 0$ and $b_n \ne 0$ for all $n$,
\[
  \frac{a_n}{b_n} \To \frac{A}{B}.
\]
Moreover the limit, when it exists, is \emph{unique}.
\end{theorem}

We will not grind out all of these -- they mirror the function case -- but the
sum law is short and shows the flavor, and uniqueness is worth a glance because
it justifies the very phrase ``the limit.''

\begin{whybox}
\textbf{Why the sum law holds.} Suppose $a_n \to A$ and $b_n \to B$. We want
$(a_n + b_n) \to A + B$. The distance to beat is
\[
  \abs{(a_n + b_n) - (A + B)} = \abs{(a_n - A) + (b_n - B)}
    \le \abs{a_n - A} + \abs{b_n - B},
\]
where the last step is the triangle inequality. Each piece on the right can be
made small. Given $\eps > 0$, use $a_n \to A$ with tolerance $\eps/2$ to get a
cutoff $N_1$ past which $\abs{a_n - A} < \eps/2$; use $b_n \to B$ with the same
tolerance to get a cutoff $N_2$ past which $\abs{b_n - B} < \eps/2$. Let $N =
\max(N_1, N_2)$. For $n \ge N$ both bounds hold at once, so
\[
  \abs{(a_n + b_n) - (A+B)} \le \abs{a_n - A} + \abs{b_n - B}
    < \tfrac\eps2 + \tfrac\eps2 = \eps .
\]
That is the whole proof. The two tricks -- split the error with the triangle
inequality, and demand $\eps/2$ from each half so they recombine to $\eps$ -- are
the bread and butter of analysis. You will use them again and again.
\end{whybox}

\begin{whybox}
\textbf{Why the limit is unique.} Suppose, for contradiction, that $a_n \to L$
and also $a_n \to L'$ with $L \ne L'$. Let $\eps = \tfrac12\abs{L - L'} > 0$,
which is the half-gap between the two candidates. Far enough out, the terms are
within $\eps$ of $L$; far enough out, they are also within $\eps$ of $L'$. Past
both cutoffs a single term $a_n$ would sit within $\eps$ of \emph{both} $L$ and
$L'$ -- but the two $\eps$-bands do not overlap, since the gap between $L$ and
$L'$ is $2\eps$. A point cannot be inside both. Contradiction. So $L = L'$: a
convergent sequence has exactly one limit.
\end{whybox}

\begin{example}[Limit laws in action]
\label{ch28:ex-rational}
Find $\displaystyle\lim_{n\to\infty} \frac{3n^2 - n + 5}{2n^2 + 7n}$.

As $n \to \infty$ both top and bottom blow up; the expression is an
$\tfrac{\infty}{\infty}$ form and we cannot read off the answer. The standard
move is to divide top and bottom by the highest power present, here $n^2$:
\[
  \frac{3n^2 - n + 5}{2n^2 + 7n}
  = \frac{3 - \dfrac{1}{n} + \dfrac{5}{n^2}}{2 + \dfrac{7}{n}} .
\]
Now each little piece is something we know. By Example
\ref{ch28:ex-oneovern}, $\tfrac1n \to 0$; multiplying that limit by itself (the
product law) gives $\tfrac{1}{n^2} \to 0$; scaling by constants and adding
(the sum and constant-multiple laws) lets us evaluate numerator and denominator
separately:
\[
  \text{numerator} \To 3 - 0 + 0 = 3,
  \qquad
  \text{denominator} \To 2 + 0 = 2 .
\]
Since the denominator's limit $2$ is nonzero, the quotient law applies and
\[
  \lim_{n\to\infty} \frac{3n^2 - n + 5}{2n^2 + 7n} = \frac{3}{2}.
\]
The leading coefficients won; everything else washed out. That is the typical
fate of a ratio of polynomials of equal degree.
\end{example}

\begin{connectionbox}
There is a clean bridge between sequence limits and the function limits of
Chapter 9. If $f$ is a function with $\lim_{x\to\infty} f(x) = L$, and $a_n =
f(n)$, then $a_n \to L$ as well -- a sequence is just the function sampled at the
integers, so if the whole continuous tail heads to $L$, the sampled dots do too.
This is what lets you import l'Hopital's rule: to find $\lim_{n\to\infty}
\tfrac{\ln n}{n}$, study the function $\tfrac{\ln x}{x}$, apply l'Hopital to get
$\tfrac{1/x}{1} \to 0$, and conclude $\tfrac{\ln n}{n} \to 0$.

The bridge runs only one way, though. The continuous limit existing forces the
sampled limit to exist, but \emph{not} conversely: $a_n = \sin(\pi n) = 0$ for
every integer $n$, so $a_n \to 0$, while $\sin(\pi x)$ oscillates forever and has
no limit as $x \to \infty$. Sampling can be calmer than the function it samples.
\end{connectionbox}

\section{The squeeze theorem for sequences}
\label{ch28:squeeze}

Some sequences resist the limit laws because they are tangled with something
that oscillates or otherwise misbehaves -- think $\dfrac{\sin n}{n}$, where the
numerator lurches between $-1$ and $1$ with no pattern you can pin down. You
cannot ``take the limit of $\sin n$'' because it has none. But you do not need
to. You only need to trap the sequence between two others that you \emph{can}
handle.

\begin{theorem}[Squeeze theorem for sequences]
\label{ch28:thm-squeeze}
Suppose that, from some index onward,
\[
  \ell_n \le a_n \le u_n,
\]
and that both bounding sequences converge to the same limit, $\ell_n \to L$ and
$u_n \to L$. Then $a_n \to L$ as well.
\end{theorem}

\begin{intuitionbox}
The picture is exactly what the name says. Two sequences -- a lower fence
$\ell_n$ and an upper fence $u_n$ -- close in on the same height $L$. The
sequence $a_n$ is caught between the fences. As the gap between the fences
crushes down to nothing, the trapped sequence has nowhere to go but $L$. It is
squeezed to its destination whether it ``wants'' to oscillate or not.
\end{intuitionbox}

\begin{figure}[H]
\centering
\begin{tikzpicture}
\begin{axis}[
  width=12cm, height=6cm,
  xlabel={$n$}, ylabel={value},
  xmin=0, xmax=21, ymin=-0.45, ymax=0.45,
  axis lines=left,
  xtick={1,5,10,15,20}, ytick={-0.4,0,0.4},
  legend pos=north east, legend style={font=\footnotesize},
  clip=false,
]
% upper fence 1/n
\addplot[targetgreen, thick, mark=none, samples=200, domain=1:20] {1/x};
\addlegendentry{$u_n=1/n$}
% lower fence -1/n
\addplot[quotientorange, thick, mark=none, samples=200, domain=1:20] {-1/x};
\addlegendentry{$\ell_n=-1/n$}
% the squeezed sequence sin(n)/n
\addplot[only marks, mark=*, mark size=1.5pt, curveblue]
  coordinates {(1,0.8415)(2,0.4546)(3,0.0470)(4,-0.1892)(5,-0.1918)
  (6,-0.0466)(7,0.0939)(8,0.1237)(9,0.0458)(10,-0.0544)(11,-0.0909)
  (12,-0.0447)(13,0.0323)(14,0.0708)(15,0.0434)(16,-0.0180)(17,-0.0566)
  (18,-0.0417)(19,0.0079)(20,0.0456)};
\addlegendentry{$a_n=\sin(n)/n$}
\addplot[densely dashed, gray] coordinates {(0,0)(21,0)};
\end{axis}
\end{tikzpicture}
\caption{The sequence $a_n = \sin(n)/n$ trapped between the fences $\ell_n =
-1/n$ and $u_n = 1/n$. Both fences converge to $0$, so the pinned sequence is
forced to $0$ -- no matter how erratically $\sin n$ jumps. (The first point
$a_1 \approx 0.84$ sits on the upper fence; thereafter the gap closes fast.)}
\label{ch28:fig-squeeze}
\end{figure}

\begin{whybox}
\textbf{Why squeezing works.} Let $\eps > 0$. Since $\ell_n \to L$ there is a
cutoff $N_1$ past which $L - \eps < \ell_n$. Since $u_n \to L$ there is a cutoff
$N_2$ past which $u_n < L + \eps$. Take $N$ to be the larger of $N_1$, $N_2$ (and
also large enough that the sandwich $\ell_n \le a_n \le u_n$ has kicked in). Then
for $n \ge N$,
\[
  L - \eps < \ell_n \le a_n \le u_n < L + \eps,
\]
so $\abs{a_n - L} < \eps$. The trapped sequence inherits the destination of its
captors. Notice we never needed a formula for $a_n$ -- only the two inequalities
and the two known limits.
\end{whybox}

\begin{example}[$\sin n / n \to 0$]
\label{ch28:ex-sinn}
For every $n$ we have $-1 \le \sin n \le 1$, and dividing by the positive number
$n$ preserves the inequalities:
\[
  -\frac1n \le \frac{\sin n}{n} \le \frac1n .
\]
Both outer sequences $\pm\tfrac1n$ converge to $0$. By the squeeze theorem,
$\dfrac{\sin n}{n} \to 0$. We tamed an untameable numerator by never touching it
-- we only bounded it. $\qed$
\end{example}

\begin{example}[A factorial beats an exponential]
\label{ch28:ex-factorial}
Show that $\displaystyle a_n = \frac{2^n}{n!} \to 0$. The factorial in the
denominator eventually overwhelms the exponential on top, but let us prove it.
Write out the factors for $n \ge 2$:
\[
  \frac{2^n}{n!}
  = \frac{2}{1}\cdot\frac{2}{2}\cdot\frac{2}{3}\cdot\frac{2}{4}\cdots\frac{2}{n}.
\]
The first two factors multiply to $2$. Every remaining factor $\tfrac{2}{k}$
(for $k \ge 3$) is below $1$, so the product of all of them is at most the single
last factor $\tfrac2n$ -- dropping factors that are below $1$ only raises the
bound. Hence
\[
  0 \le \frac{2^n}{n!} \le 2 \cdot \frac{2}{n} = \frac{4}{n}.
\]
Now $0 \to 0$ and $\tfrac4n \to 0$, so the squeeze pins $\tfrac{2^n}{n!}$ to $0$.
$\qed$
\end{example}

\section{A short table of limits worth memorizing}
\label{ch28:key-limits}

A handful of sequence limits come up so often that you should know them cold and
know \emph{why} each is true. Treat the following as a toolkit; the reasons matter
more than the entries.

\begin{center}
\renewcommand{\arraystretch}{1.5}
\begin{tabular}{@{}lll@{}}
\toprule
\textbf{Limit} & \textbf{Condition} & \textbf{Why} \\
\midrule
$\dfrac1{n^p} \to 0$ & $p > 0$ & denominator grows without bound \\
$r^n \to 0$ & $\abs{r} < 1$ & repeated shrinking factor \\
$r^n$ diverges & $\abs{r} > 1$ & repeated growing factor \\
$n^{1/n} \to 1$ & --- & the root tames the growth \\
$r^{1/n} \to 1$ & $r > 0$ & ditto, for any positive base \\
$\left(1 + \dfrac{x}{n}\right)^{\!n} \to e^x$ & $x \in \R$ & the definition of $e^x$ \\
$\dfrac{\ln n}{n^p} \to 0$ & $p > 0$ & powers beat logarithms \\
$\dfrac{n^p}{r^n} \to 0$ & $r > 1$ & exponentials beat powers \\
$\dfrac{r^n}{n!} \to 0$ & any $r$ & factorials beat exponentials \\
\bottomrule
\end{tabular}
\end{center}

\begin{happeningbox}
There is one organizing principle behind most of this table: a
\emph{hierarchy of growth}. Walk up the ladder and each rung utterly dwarfs the
one below it, eventually and forever:
\[
  \ln n \;\ll\; n^p \;\ll\; r^n \;\ll\; n! \;\ll\; n^n
  \qquad (p > 0,\ r > 1).
\]
``$\ll$'' here means the ratio of the smaller to the larger goes to $0$. A
logarithm is slower than any positive power; any positive power is slower than
any exponential with base above $1$; the factorial outruns every fixed
exponential; and $n^n$ outruns even the factorial. Whenever you face a ratio of
two such creatures, the answer is decided before you compute: the bigger one
wins, dragging the ratio to $0$ or to $\infty$ depending on who is on top.
Figure \ref{ch28:fig-ladder} shows the lower rungs racing apart on a logarithmic
scale. The
table is just this ladder, read off rung by rung. Internalize the ladder and you
can predict most sequence limits at a glance, then confirm with the squeeze
theorem or l'Hopital.
\end{happeningbox}

\begin{figure}[H]
\centering
\begin{tikzpicture}
\begin{axis}[
  width=12cm, height=6.6cm,
  xlabel={$n$}, ylabel={value (log scale)},
  ymode=log,
  xmin=1, xmax=10, ymin=0.1, ymax=5e6,
  axis lines=left,
  xtick={2,4,6,8,10},
  legend pos=north west, legend style={font=\footnotesize},
  clip=false,
]
\addplot[gray, thick, mark=*, mark size=1pt]
  coordinates {(2,0.6931)(3,1.0986)(4,1.3863)(5,1.6094)(6,1.7918)
  (7,1.9459)(8,2.0794)(9,2.1972)(10,2.3026)};
\addlegendentry{$\ln n$}
\addplot[curveblue, thick, mark=*, mark size=1pt]
  coordinates {(1,1)(2,4)(3,9)(4,16)(5,25)(6,36)(7,49)(8,64)(9,81)(10,100)};
\addlegendentry{$n^2$}
\addplot[targetgreen, thick, mark=*, mark size=1pt]
  coordinates {(1,2)(2,4)(3,8)(4,16)(5,32)(6,64)(7,128)(8,256)(9,512)(10,1024)};
\addlegendentry{$2^n$}
\addplot[trapred, thick, mark=*, mark size=1pt]
  coordinates {(1,1)(2,2)(3,6)(4,24)(5,120)(6,720)(7,5040)(8,40320)
  (9,362880)(10,3628800)};
\addlegendentry{$n!$}
\end{axis}
\end{tikzpicture}
\caption{The hierarchy of growth on a logarithmic vertical scale, where each
curve's height is its value and a higher curve grows faster. The logarithm
$\ln n$ crawls along the bottom; the power $n^2$ climbs steadily; the exponential
$2^n$ pulls ahead (note $2^n$ overtakes $n^2$ once $n \ge 5$); and the factorial
$n!$ rockets past everything. This is the ladder $\ln n \ll n^p \ll r^n \ll n!$
made visible -- and the reason a ratio of two such growth rates has its limit
decided in advance.}
\label{ch28:fig-ladder}
\end{figure}

\begin{example}[The root that tames growth: $n^{1/n} \to 1$]
\label{ch28:ex-nroot}
This one surprises people. The base $n$ races to infinity, yet the exponent
$\tfrac1n$ crashes to $0$, and the two effects very nearly cancel -- the limit is
exactly $1$. Here is a clean argument. Write $n^{1/n} = 1 + h_n$ where $h_n \ge 0$
(since $n^{1/n} \ge 1$, the excess $h_n$ is nonnegative). Raise both sides to the
$n$th power:
\[
  n = (1 + h_n)^n \ge \binom{n}{2} h_n^2 = \frac{n(n-1)}{2}\,h_n^2,
\]
where we kept just the single quadratic term from the binomial expansion -- every
other term is nonnegative, so dropping them only weakens the right side. Solve for
$h_n$:
\[
  h_n^2 \le \frac{2}{n-1}
  \quad\Longrightarrow\quad
  0 \le h_n \le \sqrt{\frac{2}{n-1}}.
\]
The right side goes to $0$, so by the squeeze $h_n \to 0$, hence $n^{1/n} = 1 +
h_n \to 1$. The trick -- keep \emph{one} well-chosen term of a binomial expansion
to get a usable inequality -- is worth filing away. $\qed$
\end{example}

\begin{example}[The compound-interest limit: $(1 + 1/n)^n \to e$]
\label{ch28:ex-euler}
You met $e$ in this guise long ago: invest a dollar at $100\%$ annual interest,
compound it $n$ times a year, and your year-end balance is $(1 + \tfrac1n)^n$. As
the compounding gets finer, the balance climbs toward a ceiling. That ceiling is
Euler's number,
\[
  \lim_{n\to\infty}\left(1 + \frac1n\right)^{\!n} = e \approx 2.71828.
\]
A few values make the climb visible: $n=1$ gives $2$; $n=2$ gives $2.25$; $n=12$
(monthly) gives about $2.613$; $n=365$ (daily) gives about $2.7146$. The sequence
increases and is bounded above (by $3$, say) -- which, by the monotone
convergence theorem of the next section, is already enough to \emph{guarantee} a
limit exists, even before we name it $e$. This is one of the cleanest examples of
a limit that defines a number rather than merely computing one.
\end{example}

\section{Monotone bounded sequences converge}
\label{ch28:monotone}

Here is the theorem that makes the real numbers worth their reputation. It says
that a sequence can be guaranteed to converge \emph{without your ever knowing
what it converges to} -- purely from two qualitative facts about its shape.

First, two pieces of vocabulary.

\begin{definition}[Monotone and bounded]
\label{ch28:def-monotone}
A sequence $(a_n)$ is \emph{increasing} if $a_n \le a_{n+1}$ for all $n$, and
\emph{decreasing} if $a_n \ge a_{n+1}$ for all $n$; it is \emph{monotone} if it
is one or the other. It is \emph{bounded above} if there is a number $M$ with
$a_n \le M$ for all $n$, \emph{bounded below} if some $m$ has $a_n \ge m$ for all
$n$, and \emph{bounded} if both.
\end{definition}

\begin{theorem}[Monotone Convergence Theorem]
\label{ch28:thm-mct}
An increasing sequence that is bounded above converges -- and its limit is the
\emph{least upper bound} (supremum) of its terms. Symmetrically, a decreasing
sequence that is bounded below converges, to the greatest lower bound of its
terms.
\end{theorem}

\begin{intuitionbox}
Picture an increasing sequence as a runner who only ever moves forward (or stands
still), but who can never cross a wall at height $M$. The runner's positions keep
rising; they are pinned below the wall. Intuitively the runner must \emph{pile
up} against some ceiling -- maybe the wall itself, maybe somewhere short of it,
but \emph{some} definite height. They cannot keep advancing forever without
either crossing the wall (forbidden) or accumulating at a limit. The limit is the
lowest height that the runner never exceeds: the least upper bound.
\end{intuitionbox}

\begin{figure}[H]
\centering
\begin{tikzpicture}
\begin{axis}[
  width=12cm, height=6cm,
  xlabel={$n$}, ylabel={$a_n$},
  xmin=0, xmax=21, ymin=0, ymax=2.4,
  axis lines=left,
  xtick={1,5,10,15,20}, ytick={1,2},
  clip=false,
]
% the ceiling / supremum (limit) = 2
\addplot[densely dashed, structurepurple, thick] coordinates {(0,2)(21,2)};
\node[structurepurple, anchor=west] at (axis cs:14.6,2.16)
  {\small wall $M=2 = \sup a_n$};
% increasing bounded sequence a_n = 2 - 2/(n+1)
\addplot[only marks, mark=*, mark size=1.6pt, curveblue, samples=20, domain=1:20]
  {2 - 2/(x+1)};
\end{axis}
\end{tikzpicture}
\caption{An increasing sequence ($a_n = 2 - \tfrac{2}{n+1}$) bounded above by the
wall $M = 2$. The terms climb but never cross; they pile up against the least
upper bound, which is exactly the limit. Monotone plus bounded forces a
destination.}
\label{ch28:fig-mct}
\end{figure}

\begin{whybox}
\textbf{Why monotone + bounded forces convergence.} This is where the real
numbers do their secret work. The natural numbers, the rationals -- they are not
enough. The decisive property is \emph{completeness}: every nonempty set of reals
that is bounded above has a \emph{least} upper bound, a supremum, denoted $\sup$.
There are no ``gaps'' for a least upper bound to fall into. (The rationals lack
this: the rationals less than $\sqrt2$ are bounded above but have no rational
least upper bound, because $\sqrt2$ is missing from $\Q$.)

Now take an increasing sequence $(a_n)$ bounded above. Let $L = \sup\{a_n\}$, the
least upper bound of the set of its terms; completeness guarantees this number
exists. We claim $a_n \to L$. Let $\eps > 0$. Because $L$ is the \emph{least}
upper bound, $L - \eps$ is \emph{not} an upper bound -- something smaller than $L$
fails to bound the terms -- so some term $a_N$ exceeds $L - \eps$. But the
sequence is increasing, so for every $n \ge N$ we have $a_n \ge a_N > L - \eps$.
And every term is $\le L$ since $L$ bounds them all. Therefore, for $n \ge N$,
\[
  L - \eps < a_n \le L,
  \qquad\text{so}\qquad
  \abs{a_n - L} < \eps .
\]
That is convergence to $L$. The engine is completeness; monotonicity is what lets
a single term $a_N$ vouch for all that follow it. $\qed$
\end{whybox}

\begin{aestheticbox}
Sit with what this theorem buys you. To prove a sequence converges directly from
Definition \ref{ch28:def-limit}, you must already \emph{know} the limit $L$ --
you cannot check $\abs{a_n - L} < \eps$ without an $L$ to subtract. But countless
sequences in mathematics are defined in ways that hide their limit completely:
the partial sums of a series, the iterates of Newton's method, the successive
decimal truncations of an unknown real number. The Monotone Convergence Theorem
hands you the limit's \emph{existence} for free, in exchange for two facts you
can often check by hand -- ``it goes up'' and ``it stays below something.'' You
get the destination guaranteed without buying a map. This is why so many
existence proofs in analysis end with the phrase ``monotone and bounded, hence
convergent.''
\end{aestheticbox}

\begin{connectionbox}
This theorem is the precise reason an infinite decimal means anything. The string
$0.333\ldots$ is the sequence of truncations $0.3,\ 0.33,\ 0.333,\ \dots$. Each
truncation is bigger than the last (increasing), and every one is below $0.4$
(bounded above). Monotone plus bounded: the sequence converges, to a genuine real
number -- the one we cheerfully call $\tfrac13$. Without the Monotone Convergence
Theorem, ``the number $0.333\ldots$'' would be a string of digits with no
guaranteed value. The same logic certifies \emph{every} infinite decimal as a
real number. Chapter 29 builds infinite series on exactly this foundation:
a series converges precisely when its sequence of partial sums does.
\end{connectionbox}

\section{Recursive sequences and fixed points}
\label{ch28:recursion}

Many of the most important sequences are not given by a formula in $n$ but by a
\emph{rule for getting the next term from the current one}:
\[
  a_{n+1} = f(a_n),
\]
starting from some seed $a_1$. Each term is the previous term fed through the
same function $f$. Newton's method is of this form; so is every numerical scheme
that ``iterates to a solution.'' The natural questions are: does the sequence
settle down, and if so, where?

\begin{happeningbox}
\textbf{What is really happening: iteration and fixed points.} Suppose the
sequence $a_{n+1} = f(a_n)$ \emph{does} converge, to some limit $L$. Then both
$a_n \to L$ and $a_{n+1} \to L$ (a sequence and its own shifted copy share a
limit -- shifting by one place cannot change where a list is headed). If $f$ is
continuous, taking the limit of both sides of $a_{n+1} = f(a_n)$ gives
\[
  L = f(L).
\]
A point satisfying $L = f(L)$ is called a \emph{fixed point} of $f$: feeding it
through $f$ returns it unchanged. So \emph{the only possible limits of the
iteration are the fixed points of $f$}. This is the master key. The recursion can
only come to rest where $f$ leaves a point alone. To find the candidate limits,
you do not analyze the sequence at all -- you solve the equation $L = f(L)$.
\end{happeningbox}

\begin{figure}[H]
\centering
\begin{tikzpicture}[scale=2.6]
% axes
\draw[->] (0,0) -- (1.15,0) node[right] {$x$};
\draw[->] (0,0) -- (0,1.15) node[above] {$y$};
% line y=x
\draw[densely dashed, gray] (0,0) -- (1.1,1.1) node[right, black] {\small $y=x$};
% curve y = f(x) = sqrt(x) (concave, crosses y=x at 0 and 1)
\draw[curveblue, thick] plot[domain=0.02:1.08, samples=60] (\x, {sqrt(\x)})
  node[right, black] {\small $y=f(x)$};
% fixed point at (1,1)
\fill[trapred] (1,1) circle (0.013);
\node[trapred, anchor=north west] at (1,1) {\small $L$};
% cobweb starting at a1 = 0.15
\draw[targetgreen, thick] (0.15,0) -- (0.15,0.387); % up to curve
\draw[targetgreen, thick] (0.15,0.387) -- (0.387,0.387); % across to line
\draw[targetgreen, thick] (0.387,0.387) -- (0.387,0.622);
\draw[targetgreen, thick] (0.387,0.622) -- (0.622,0.622);
\draw[targetgreen, thick] (0.622,0.622) -- (0.622,0.789);
\draw[targetgreen, thick] (0.622,0.789) -- (0.789,0.789);
\draw[targetgreen, thick] (0.789,0.789) -- (0.789,0.888);
\draw[targetgreen, thick] (0.789,0.888) -- (0.888,0.888);
\node[anchor=north] at (0.15,0) {\small $a_1$};
\end{tikzpicture}
\caption{A \emph{cobweb diagram} for the iteration $a_{n+1} = f(a_n)$ with
$f(x)=\sqrt{x}$. Starting from $a_1$, go vertically to the curve (apply $f$),
then horizontally to the line $y=x$ (copy the output back to the input axis), and
repeat. The staircase spirals into the fixed point $L=1$, where the curve meets
the diagonal. Fixed points are exactly the crossings of $y=f(x)$ and $y=x$.}
\label{ch28:fig-cobweb}
\end{figure}

\begin{intuitionbox}
\textbf{Reading a cobweb.} The line $y = x$ is the ``mirror'' that copies an
output back into an input. A fixed point is where the curve $y = f(x)$ crosses
this mirror -- there, output equals input, so the iteration stops moving. Whether
the staircase is drawn \emph{toward} a fixed point (it attracts) or \emph{away}
from it (it repels) is governed by how steep $f$ is there: a shallow curve
(slope of magnitude below $1$) pulls nearby points in; a steep one (slope above
$1$) pushes them away. The next box makes that precise.
\end{intuitionbox}

\begin{example}[A recursion that converges: $a_{n+1} = \sqrt{2 + a_n}$]
\label{ch28:ex-nested-radical}
Define $a_1 = \sqrt2$ and $a_{n+1} = \sqrt{2 + a_n}$. The terms are the nested
radicals
\[
  \sqrt2,\quad \sqrt{2 + \sqrt2},\quad \sqrt{2 + \sqrt{2 + \sqrt2}},\quad \dots
\]
numerically $1.414,\ 1.848,\ 1.962,\ 1.990,\ 1.998,\ \dots$, visibly climbing
toward $2$. Let us prove they converge and find the limit, the right way.

\textbf{Step 1: find the candidate limit.} If $a_n \to L$, then $L$ is a fixed
point of $f(x) = \sqrt{2 + x}$, so $L = \sqrt{2 + L}$. Square: $L^2 = 2 + L$, i.e.
$L^2 - L - 2 = 0$, i.e. $(L-2)(L+1) = 0$. The roots are $L = 2$ and $L = -1$.
Since every term is positive, the only viable candidate is $L = 2$.

\textbf{Step 2: show it is bounded above (by $2$).} Use induction. $a_1 =
\sqrt2 < 2$. If $a_n < 2$, then $a_{n+1} = \sqrt{2 + a_n} < \sqrt{2 + 2} =
\sqrt4 = 2$. So $a_n < 2$ for all $n$.

\textbf{Step 3: show it is increasing.} We check $a_{n+1} > a_n$, i.e. $\sqrt{2 +
a_n} > a_n$. Squaring (both sides positive), this is $2 + a_n > a_n^2$, i.e.
$a_n^2 - a_n - 2 < 0$, i.e. $(a_n - 2)(a_n + 1) < 0$. Since $0 < a_n < 2$ from
Step 2, the first factor is negative and the second positive, so the product is
negative. The inequality holds; the sequence increases.

\textbf{Step 4: conclude.} The sequence is increasing and bounded above, so by
the Monotone Convergence Theorem (\ref{ch28:thm-mct}) it converges. By Step 1 the
limit can only be $2$. Therefore
\[
  \sqrt{2 + \sqrt{2 + \sqrt{2 + \cdots}}} = 2. \qquad\qed
\]
Notice the division of labor: Step 1 (fixed point) tells us \emph{where} the
limit must be; Steps 2--3 (monotone, bounded) tell us a limit \emph{exists at
all}. Neither alone is enough -- the fixed-point equation also offered $-1$, a
phantom the boundedness argument quietly rules out.
\end{example}

\begin{whybox}
\textbf{Why slope decides attraction.} Suppose $L$ is a fixed point, $f$ is
differentiable, and $\abs{f'(L)} < 1$. Near $L$ the map moves a point at distance
$d$ from $L$ to a point at distance roughly $\abs{f'(L)}\,d$ -- closer, because
the multiplier is below $1$. Each step shrinks the distance by a factor under
$1$, so the errors form a sequence dominated by a geometric one tending to $0$:
the iteration is drawn in, and we call $L$ \emph{attracting} (stable). If instead
$\abs{f'(L)} > 1$, each step \emph{enlarges} the distance and the iteration is
flung away ($L$ is \emph{repelling}). At the borderline $\abs{f'(L)} = 1$ the
linear term decides nothing and you must look closer.

In Example \ref{ch28:ex-nested-radical}, $f(x) = \sqrt{2+x}$ has $f'(x) =
\tfrac{1}{2\sqrt{2+x}}$, so $f'(2) = \tfrac{1}{2\cdot 2} = \tfrac14 < 1$: the
fixed point $2$ attracts, consistent with the staircase marching in. The phantom
root $-1$ is not even in the domain we care about, so it never competes.
\end{whybox}

\begin{example}[A recursion that does \emph{not} settle]
\label{ch28:ex-divergent-recursion}
Not every recursion converges. Take $a_{n+1} = 2a_n - 1$ with seed $a_1 = 2$. The
fixed-point equation $L = 2L - 1$ gives $L = 1$. But $f(x) = 2x - 1$ has $f'(x) =
2 > 1$ everywhere, so the fixed point $1$ repels. Starting at $a_1 = 2$ (distance
$1$ above the fixed point) the terms are
\[
  2,\ 3,\ 5,\ 9,\ 17,\ 33,\ \dots,
\]
each step doubling the distance from $1$. The sequence diverges to $+\infty$. The
fixed point exists, but the iteration runs from it rather than to it. This is the
cautionary twin of the previous example: solving $L = f(L)$ tells you the
\emph{only possible} limit, never that a limit is \emph{achieved}. Convergence
must be earned separately -- by monotone-bounded, by the slope test, or by direct
estimate.
\end{example}

\begin{historybox}
The fixed-point idea is one of the great unifying threads of modern mathematics.
In 1922 Stefan Banach proved his \emph{contraction mapping theorem}: if a map $f$
shrinks distances by a fixed factor below $1$ -- $\abs{f(x) - f(y)} \le k\,
\abs{x-y}$ with $k < 1$ -- then iterating from \emph{any} starting point produces
a sequence converging to the one and only fixed point. The slope condition
$\abs{f'(L)} < 1$ you met above is the local, calculus-flavored shadow of
Banach's global statement.

This is not abstraction for its own sake. Banach's theorem is the engine behind
the existence-and-uniqueness proof for differential equations (Picard
iteration, Chapter 37), behind Newton's method, behind Google's PageRank, and
behind countless numerical solvers. In each case a hard problem is rephrased as
``find the fixed point of this map,'' and a sequence of approximations is set
loose to crawl toward it. The humble recursion $a_{n+1} = f(a_n)$ of this section
is the seed of all of it.
\end{historybox}

\begin{applicationbox}
\textbf{Babylonian square roots.} To compute $\sqrt{S}$, iterate
\[
  a_{n+1} = \frac12\left(a_n + \frac{S}{a_n}\right),
\]
which is exactly Newton's method applied to $x^2 - S = 0$. Its fixed points solve
$L = \tfrac12(L + S/L)$, i.e. $2L = L + S/L$, i.e. $L^2 = S$, i.e. $L = \sqrt S$.
For $S = 2$ starting at $a_1 = 1$:
\[
  1,\quad 1.5,\quad 1.41667,\quad 1.41422,\quad 1.41421356\ldots,
\]
and the correct figures roughly \emph{double} each step -- a hallmark of Newton's
method near a simple root. The Babylonians ran this recursion on clay tablets
nearly four thousand years ago. They had no theory of limits, no $\eps$, no $N$ --
only the observation that the iteration ``stopped changing.'' What this chapter
supplies is the why: $\tfrac12(x + S/x)$ has an attracting fixed point at
$\sqrt S$, so the clay tablets were converging all along.
\end{applicationbox}

\begin{labbox}
Pick a recursion and run it by hand for five steps, then explain its fate.
\begin{enumerate}
  \item $a_1 = 0$, $a_{n+1} = \cos(a_n)$ (radians). Compute $a_2,\dots,a_6$ on a
        calculator. The terms hover around a single number -- the
        \emph{Dottie number}. What equation does that number satisfy? (Answer:
        the fixed point $L = \cos L \approx 0.739$.) Is $\abs{f'(L)} < 1$?
  \item $a_1 = 1$, $a_{n+1} = 1 + \tfrac1{a_n}$. Compute several terms. They
        approach the golden ratio $\varphi$. Find the fixed-point equation $L =
        1 + \tfrac1L$ and solve it. Which root do the terms choose, and why?
\end{enumerate}
\end{labbox}

\section{From sequences to series: a preview}
\label{ch28:to-series}

Everything in this chapter was built to support one move, the move that occupies
the next several chapters: making sense of an \emph{infinite sum}.

You cannot literally add infinitely many numbers -- addition is an operation on
two inputs, repeated finitely. So what could
\[
  \tfrac12 + \tfrac14 + \tfrac18 + \tfrac1{16} + \cdots
\]
possibly \emph{mean}? The only sensible answer is a sequence in
disguise. Form the \emph{partial sums}, the running totals:
\[
  s_1 = \tfrac12,\quad
  s_2 = \tfrac12 + \tfrac14 = \tfrac34,\quad
  s_3 = \tfrac78,\quad
  s_4 = \tfrac{15}{16},\quad \dots,\quad
  s_n = 1 - \tfrac1{2^n}.
\]
This is a perfectly ordinary sequence $(s_n)$. It is increasing (each new term
adds a positive amount) and bounded above (by $1$). By the Monotone Convergence
Theorem it converges -- and clearly to $1$. \emph{That} is what we mean when we
write the infinite sum and declare it equal to $1$: the sequence of partial sums
converges to $1$.

\begin{figure}[H]
\centering
\begin{tikzpicture}
\begin{axis}[
  width=12cm, height=5.6cm,
  xlabel={$n$}, ylabel={$s_n$},
  xmin=0, xmax=11, ymin=0, ymax=1.2,
  axis lines=left,
  xtick={1,...,10}, ytick={0.5,1},
  clip=false,
]
\addplot[densely dashed, structurepurple, thick] coordinates {(0,1)(11,1)};
\node[structurepurple, anchor=west] at (axis cs:7.4,1.08) {\small limit $=1$};
\addplot[only marks, mark=*, mark size=1.7pt, curveblue]
  coordinates {(1,0.5)(2,0.75)(3,0.875)(4,0.9375)(5,0.96875)(6,0.984375)
  (7,0.9921875)(8,0.99609375)(9,0.998046875)(10,0.9990234375)};
\end{axis}
\end{tikzpicture}
\caption{Partial sums $s_n = 1 - 2^{-n}$ of the series $\tfrac12 + \tfrac14 +
\tfrac18 + \cdots$. The running totals climb, halving the remaining gap to $1$ at
each step: increasing and bounded above, hence convergent. The infinite sum is
\emph{defined} to be this limit, $1$.}
\label{ch28:fig-partialsums}
\end{figure}

\begin{aestheticbox}
This is the whole strategy of the infinite, in one sentence: \textbf{an infinite
process is defined to be the limit of its finite stages.} An infinite sum is the
limit of partial sums. An infinite product is the limit of partial products. A
real number is the limit of its decimal truncations. A definite integral is the
limit of Riemann sums. In every case the actual infinity is never performed --
it is replaced by a sequence of finite, computable approximations, and
the ``answer'' is wherever that sequence is \emph{headed}. The theory of
sequences you have just built is therefore not a side topic. It is the
foundation on which the entire calculus of the infinite stands. Master where a
list of numbers goes, and you have mastered, in advance, where every infinite
process in this book goes.
\end{aestheticbox}

\chapterline

\section*{Exercises}
\addcontentsline{toc}{section}{Exercises}

\begin{exercise}
\label{ch28:exr-terms}
Write out the first five terms of each sequence and predict, with a one-line
reason, where it is headed.
\begin{enumerate}[label=(\alph*)]
  \item $a_n = \dfrac{2n - 1}{n + 3}$
  \item $a_n = \dfrac{(-1)^{n+1}}{2^n}$
  \item $a_n = \dfrac{n!}{n^n}$
  \item $a_n = \cos(\pi n)$
\end{enumerate}
\end{exercise}

\begin{exercise}
\label{ch28:exr-epsilonN}
Use the epsilon-$N$ definition (\ref{ch28:def-limit}) directly to prove each
limit. In each case, given $\eps > 0$, exhibit an explicit $N$ and verify the
implication.
\begin{enumerate}[label=(\alph*)]
  \item $\displaystyle\lim_{n\to\infty} \frac{1}{n^2} = 0$.
  \item $\displaystyle\lim_{n\to\infty} \frac{n}{n+1} = 1$.
        \emph{Hint: simplify $\abs{\tfrac{n}{n+1} - 1}$ first.}
  \item $\displaystyle\lim_{n\to\infty} \frac{3n + 2}{n} = 3$.
\end{enumerate}
\end{exercise}

\begin{exercise}
\label{ch28:exr-laws}
Evaluate each limit, naming the limit laws or table entries
(Section \ref{ch28:key-limits}) you use.
\begin{enumerate}[label=(\alph*)]
  \item $\displaystyle\lim_{n\to\infty} \frac{4n^3 - 2n + 1}{5n^3 + n^2}$
  \item $\displaystyle\lim_{n\to\infty} \frac{\sqrt{n^2 + 1}}{n}$
  \item $\displaystyle\lim_{n\to\infty} \left(\sqrt{n+1} - \sqrt{n}\right)$
        \emph{Hint: multiply by the conjugate.}
  \item $\displaystyle\lim_{n\to\infty} \frac{2^n + 3^n}{3^n}$
\end{enumerate}
\end{exercise}

\begin{exercise}
\label{ch28:exr-squeeze}
Use the squeeze theorem (\ref{ch28:thm-squeeze}) to evaluate each limit. State
your two bounding sequences explicitly.
\begin{enumerate}[label=(\alph*)]
  \item $\displaystyle\lim_{n\to\infty} \frac{\cos(n^2)}{n}$
  \item $\displaystyle\lim_{n\to\infty} \frac{n + \sin n}{n^2}$
  \item $\displaystyle\lim_{n\to\infty} \frac{3^n}{n!}$ \emph{(adapt Example
        \ref{ch28:ex-factorial})}
  \item $\displaystyle\lim_{n\to\infty} \sqrt[n]{2^n + 3^n}$
        \emph{Hint: bound below by $3$ and above by $3\sqrt[n]{2}$.}
\end{enumerate}
\end{exercise}

\begin{exercise}
\label{ch28:exr-oscillate}
Explain carefully, using the definition of a limit, why the sequence $a_n =
(-1)^n \cdot \dfrac{n}{n+1}$ diverges. (What two values do the even and odd
terms approach, and why does that forbid a single limit?)
\end{exercise}

\begin{exercise}
\label{ch28:exr-mct-recursion}
Let $a_1 = 1$ and $a_{n+1} = \sqrt{6 + a_n}$.
\begin{enumerate}[label=(\alph*)]
  \item Find the candidate limit by solving the fixed-point equation $L =
        \sqrt{6 + L}$; discard the inadmissible root.
  \item Prove by induction that $a_n < 3$ for all $n$ (bounded above).
  \item Prove the sequence is increasing.
  \item Conclude that it converges, and state the limit.
\end{enumerate}
\end{exercise}

\begin{exercise}
\label{ch28:exr-decreasing}
Let $a_1 = 3$ and $a_{n+1} = \dfrac{a_n + 4}{2}$.
\begin{enumerate}[label=(\alph*)]
  \item Compute the first four terms. Is the sequence increasing or decreasing?
  \item Find the fixed point of $f(x) = \tfrac{x+4}{2}$ and compute $f'$; is the
        fixed point attracting or repelling?
  \item Show the sequence is bounded above by the fixed point, then conclude it
        converges (which form of the Monotone Convergence Theorem applies here?).
\end{enumerate}
\end{exercise}

\begin{exercise}
\label{ch28:exr-cobweb}
Consider the logistic-style map $a_{n+1} = \tfrac32\,a_n(1 - a_n)$ with seed
$a_1 = 0.2$.
\begin{enumerate}[label=(\alph*)]
  \item Find the two fixed points of $f(x) = \tfrac32 x(1-x)$.
  \item Compute $f'$ at each fixed point and classify each as attracting or
        repelling.
  \item Compute the first five terms numerically and say which fixed point the
        iteration is approaching, consistent with part (b).
\end{enumerate}
\end{exercise}

\begin{exercise}
\label{ch28:exr-babylonian}
Apply the Babylonian iteration $a_{n+1} = \tfrac12(a_n + \tfrac{5}{a_n})$ with
$a_1 = 2$ to approximate $\sqrt5$.
\begin{enumerate}[label=(\alph*)]
  \item Compute $a_2, a_3, a_4$ to six decimal places.
  \item How many correct digits does each step produce? Comment on the doubling
        you observe and connect it to the discussion in the application box.
\end{enumerate}
\end{exercise}

\begin{exercise}
\label{ch28:exr-e}
\begin{enumerate}[label=(\alph*)]
  \item Using a calculator, tabulate $(1 + \tfrac1n)^n$ for $n = 1, 10, 100,
        1000, 10{,}000$ and watch it climb toward $e$.
  \item The sequence is increasing and bounded above. Which theorem of this
        chapter guarantees, before any numerical evidence, that the limit
        \emph{exists}? Explain in one or two sentences what role completeness of
        the real numbers plays.
\end{enumerate}
\end{exercise}

\begin{exercise}
\label{ch28:exr-growth-ladder}
Without computing exact values, use the growth hierarchy $\ln n \ll n^p \ll r^n
\ll n! \ll n^n$ (Section \ref{ch28:key-limits}) to decide whether each limit is
$0$ or $+\infty$, and justify in one line.
\begin{enumerate}[label=(\alph*)]
  \item $\displaystyle\lim_{n\to\infty} \frac{n^{10}}{2^n}$
  \item $\displaystyle\lim_{n\to\infty} \frac{n!}{100^n}$
  \item $\displaystyle\lim_{n\to\infty} \frac{(\ln n)^3}{\sqrt n}$
  \item $\displaystyle\lim_{n\to\infty} \frac{n^n}{n!}$
\end{enumerate}
\end{exercise}

\begin{exercise}
\label{ch28:exr-decimals}
A real number is the limit of its decimal truncations. Let $a_n$ be $\pi$
truncated to $n$ decimal places: $a_1 = 3.1,\ a_2 = 3.14,\ a_3 = 3.141, \dots$.
\begin{enumerate}[label=(\alph*)]
  \item Argue that $(a_n)$ is increasing and bounded above (by $3.2$, say).
  \item Which theorem certifies that this sequence has a limit -- so that ``the
        number $\pi$'' is a well-defined real number rather than just an endless
        string of digits? Relate your answer to the connection box on infinite
        decimals.
\end{enumerate}
\end{exercise}
